\documentclass[a4paper]{article}

\usepackage[fontset=fandol]{ctex}
\usepackage{amsmath, amssymb}
\usepackage{graphicx}
\usepackage[margin=2.45cm]{geometry}
\usepackage[most]{tcolorbox}
\usepackage{hyperref}
\usepackage{booktabs}
\usepackage{float}
\usepackage{array}
\usepackage{xcolor}
\usepackage{enumitem}
\usepackage{caption}

\setlength{\parindent}{2em}
\setlength{\parskip}{0.35em}
\setlength{\emergencystretch}{3em}
\linespread{1.06}
\sloppy
\raggedbottom

\newcolumntype{L}[1]{>{\raggedright\arraybackslash}p{#1}}
\newcolumntype{C}[1]{>{\centering\arraybackslash}p{#1}}

\DeclareMathOperator*{\argmin}{arg\,min}
\DeclareMathOperator*{\argmax}{arg\,max}
\DeclareMathOperator{\softmax}{softmax}
\DeclareMathOperator{\Cat}{Cat}
\DeclareMathOperator{\E}{\mathbb{E}}
\DeclareMathOperator{\KL}{KL}

\hypersetup{
    colorlinks=true,
    linkcolor=blue!60!black,
    urlcolor=blue!60!black
}

\setlist[itemize]{leftmargin=2.2em,itemsep=0.22em,topsep=0.25em}
\setlist[enumerate]{leftmargin=2.4em,itemsep=0.22em,topsep=0.25em}
\setlist[description]{leftmargin=3.0em,itemsep=0.18em,topsep=0.25em}

\newtcolorbox{knowledgebox}[1]{
    enhanced, colback=blue!5!white, colframe=blue!75!black, colbacktitle=blue!75!black,
    coltitle=white, fonttitle=\bfseries, title={#1},
    attach boxed title to top left={yshift=-2mm, xshift=2mm},
    boxrule=1pt, sharp corners
}

\newtcolorbox{importantbox}[1]{
    enhanced, colback=yellow!10!white, colframe=yellow!80!black, colbacktitle=yellow!80!black,
    coltitle=black, fonttitle=\bfseries, title={#1}, sharp corners
}

\newtcolorbox{warningbox}[1]{
    enhanced, colback=red!5!white, colframe=red!75!black, colbacktitle=red!75!black,
    coltitle=white, fonttitle=\bfseries, title={#1}, sharp corners
}

\newcommand{\notetitle}{自然语言处理上的对抗式攻击（下）：生成式攻击与防御}
\newcommand{\notesubtitle}{Auto-Encoder Attacks, Gumbel-Softmax, Robust Training, ASCC, and Detection}
\newcommand{\noteauthors}{基于李宏毅老师《机器学习 2025》课程内容整理}
\newcommand{\notedate}{\today}
\newcommand{\videochannel}{李宏毅深度学习课堂（Bilibili）}
\newcommand{\videoduration}{约 1 小时 10 分 19 秒}
\newcommand{\videourl}{https://www.bilibili.com/video/BV1TAtwzTE1S?p=42}

\newcommand{\framefig}[4]{%
\begin{figure}[H]
    \centering
    \includegraphics[width=#1\textwidth]{#2}
    \caption{#3\protect\footnotemark}
\end{figure}
\footnotetext{视频画面时间区间：#4。}
}

\begin{document}
\pagestyle{empty}

\begin{center}
    \vspace*{1.0cm}
    {\Huge\bfseries \notetitle\par}
    \vspace{0.35cm}
    {\LARGE \notesubtitle\par}
    \vspace{0.75cm}
    \includegraphics[width=0.62\textwidth]{video.jpg}\par
    \vspace{0.75cm}
    {\large \noteauthors\par}
    {\large \notedate\par}
    \vspace{0.75cm}
    \begin{tcolorbox}[width=0.86\textwidth,center,sharp corners,
                      colback=black!3!white,colframe=black!50!white]
        \centering
        \textbf{视频信息}\\[3pt]
        频道：\videochannel \\
        时长：\videoduration \\
        链接：\href{\videourl}{\videourl}
    \end{tcolorbox}
\end{center}

\newpage
\tableofcontents
\newpage
\pagestyle{plain}

% =====================================================================
\section{从词替换到生成式攻击}
% =====================================================================

上一讲主要围绕 synonym substitution attack：攻击者先找重要词，再从同义词或语言模型候选中挑替换词。P42 继续讨论 evasion attack，但重点从“逐词替换”转向“训练一个生成器”。如果有一个 generator 能自动把原句 $x$ 改成 $x'$，同时让 $x'$ 保持语义、语法自然并欺骗分类器，那么攻击就可以被写成一个可训练的生成任务。

这条路线的动机很直接：手写搜索规则会遇到大量启发式选择，例如候选集合怎么建、先替换哪个词、约束阈值设多少。生成式方法希望让模型自己学会产生扰动句子，再通过分类器 loss、重建 loss 和语义 loss 把它限制在合理范围内。

\framefig{0.90}{figures/fig01_autoencoder_generator_goal.jpg}
{Auto-encoder attack 的基本目标：训练 generator 产生 adversarial sample，使固定的 text classifier 预测错误。}
{00:02:15--00:02:30}

\subsection{Generator、classifier 与攻击目标}

设文本分类器为 $C$，原始输入为 $x$，真实标签为 $y$。生成器或 auto-encoder 记为 $G_\theta$，它把 $x$ 改写成：
\[
x' = G_\theta(x).
\]

攻击者希望 $x'$ 同时满足两件事：

\begin{enumerate}
    \item 对人类而言，$x'$ 与 $x$ 足够接近，语义没有被偷换。
    \item 对模型而言，$C(x')$ 不再输出 $y$，或输出攻击者指定的目标类别 $y_{\mathrm{tar}}$。
\end{enumerate}

对 untargeted attack，可以让 adversarial loss 鼓励分类器降低真实标签概率：
\[
L_{\mathrm{adv}}(x') = \log p_C(y\mid x')
\quad\text{或}\quad
-L_C(C(x'),y),
\]
具体符号取决于实现是最小化还是最大化。直观上，生成器要让分类器对正确类别越来越没有信心。对 targeted attack，则可以改成最小化：
\[
L_{\mathrm{tar}}(x')=-\log p_C(y_{\mathrm{tar}}\mid x').
\]

\subsection{重建、语义与攻击 loss}

生成器不能只追求让分类器出错，否则它可能输出完全无关的句子。课程用 reconstruction、similarity 和 adversarial loss 共同约束生成过程。

\framefig{0.90}{figures/fig02_autoencoder_attack_losses.jpg}
{攻击阶段的 loss 通常同时包含重建项、语义相似项和诱导分类器出错的 adversarial loss。}
{00:05:30--00:05:45}

可以把训练目标概括成：
\[
\min_\theta
\lambda_{\mathrm{rec}}L_{\mathrm{rec}}(x,G_\theta(x))
+\lambda_{\mathrm{sim}}L_{\mathrm{sim}}(x,G_\theta(x))
+\lambda_{\mathrm{adv}}L_{\mathrm{adv}}(G_\theta(x)).
\]

其中：

\begin{description}
    \item[$L_{\mathrm{rec}}$] 要求 auto-encoder 不要离原句太远，可理解为重建原句的语言建模损失。
    \item[$L_{\mathrm{sim}}$] 要求生成句与原句在 embedding 或句向量空间中接近。
    \item[$L_{\mathrm{adv}}$] 要求生成句让固定分类器犯错。
    \item[$\lambda$] 控制三种目标的相对权重。
\end{description}

这些项之间天然冲突。若 $L_{\mathrm{rec}}$ 太强，生成器只会复读原句，攻击失败；若 $L_{\mathrm{adv}}$ 太强，生成器可能生成语义崩坏的句子；若 $L_{\mathrm{sim}}$ 设计得太弱，成功率会被低质量样本夸大。

\subsection{攻击与防御的交替视角}

同一个 auto-encoder 生成的样本也可以用于 defense。攻击阶段固定分类器、更新生成器；防御阶段可以固定或使用生成器产生的 adversarial examples，再更新分类器，让分类器学会把它们判回正确标签。

\framefig{0.90}{figures/fig03_autoencoder_defense_loss.jpg}
{防御阶段会把生成的 adversarial sample 喂回分类器，使分类器在这些扰动输入上仍预测正确。}
{00:08:45--00:09:00}

这对应到 adversarial training 的基本形式：
\[
\min_\phi
\E_{(x,y)\sim\mathcal{D}}
\left[
L_C(C_\phi(x),y)
+\beta L_C(C_\phi(G_\theta(x)),y)
\right].
\]

第一项维持 clean input 的表现，第二项让模型在生成器产生的 $x'$ 上也预测正确。这里的 $G_\theta$ 可以是攻击者训练出来的，也可以是防御者内部用来生成训练样本的攻击模块。

\begin{knowledgebox}{生成式攻击的核心张力}
    生成器越自由，越容易找到让分类器出错的输入；生成器越受约束，越能保证样本自然、语义一致，但攻击越困难。NLP adversarial generation 的难点，就是把“可攻击性”和“可读性”同时放进训练目标。
\end{knowledgebox}

\subsection{本章小结}

Auto-encoder attack 把文本对抗样本生成看成一个训练问题：generator 负责改写句子，classifier 提供攻击信号，reconstruction 和 similarity loss 保持文本质量。它比纯词替换搜索更“模型化”，但也引出了一个关键技术问题：文字是离散的，采样和 argmax 不可微，梯度如何从分类器传回 generator？

% =====================================================================
\section{离散采样不可微与 Gumbel-Softmax}
% =====================================================================

图像生成器输出连续像素，分类器 loss 可以直接反向传播到生成器。文本生成器不同：它通常输出词表上的概率分布，然后通过 sampling 或 argmax 选出一个 token，再查 embedding table 送入分类器。这个离散选择步骤不可微，导致分类器的梯度无法直接穿过采样步骤更新 auto-encoder。

\framefig{0.90}{figures/fig04_non_differentiable_sampling.jpg}
{Auto-encoder 输出词表分布后若直接采样 token，再做 embedding lookup，采样位置会阻断梯度。}
{00:14:00--00:14:15}

\subsection{为什么 argmax 会阻断梯度}

假设生成器在某个位置输出 categorical distribution：
\[
\pi=(\pi_1,\pi_2,\dots,\pi_V),\qquad \sum_{i=1}^{V}\pi_i=1.
\]

离散生成通常取：
\[
z=\mathrm{onehot}(\argmax_i \pi_i).
\]

问题在于 $\argmax$ 对输入概率的微小变化大多没有反应；一旦最大项没变，输出 one-hot 完全不变；最大项切换时又是突变。这样的函数几乎处处梯度为零，切换点不可导。若后面的 classifier loss 为 $L(C(Ez),y)$，梯度无法稳定回到 $\pi$。

这也是 NLP 生成式攻击比图像生成式攻击麻烦的地方：模型真正想优化的是离散 token，但梯度优化喜欢连续空间。

\subsection{Gumbel-Max trick}

Gumbel-Max trick 提供了一种从 categorical distribution 中采样的写法。若 $g_i$ 是独立同分布的 Gumbel noise：
\[
g_i=-\log(-\log u_i),\qquad u_i\sim \mathrm{Uniform}(0,1),
\]
则可以用：
\[
z=\mathrm{onehot}\left(\argmax_i(\log \pi_i+g_i)\right)
\]
从 $\Cat(\pi)$ 中采样。

\framefig{0.90}{figures/fig05_gumbel_softmax_intro.jpg}
{Gumbel-Softmax reparameterization trick 从 Gumbel-Max 出发，把离散采样改写成可近似微分的形式。}
{00:14:15--00:14:30}

Gumbel-Max 本身仍包含 $\argmax$，所以还不可微；但它把随机性显式写成 $g_i$，为后续的 softmax relaxation 铺路。

\subsection{Softmax temperature relaxation}

Gumbel-Softmax 用带温度 $T$ 的 softmax 近似 $\argmax$：
\[
y_i=
\frac{\exp((\log \pi_i+g_i)/T)}
{\sum_{j=1}^{V}\exp((\log \pi_j+g_j)/T)}.
\]

\framefig{0.90}{figures/fig06_gumbel_argmax_reparameterization.jpg}
{Gumbel-Softmax 用 softmax with temperature scaling 近似 argmax，使 categorical sampling 有连续替代。}
{00:19:45--00:20:00}

当 $T$ 很小时，最大的那一项会接近 1，其余接近 0，$y$ 接近 one-hot；当 $T$ 较大时，分布更平滑，梯度更稳定但离离散 token 更远。温度 $T$ 因此控制了“可微性”和“离散性”的折中。

\framefig{0.90}{figures/fig07_gumbel_softmax_solution.jpg}
{Gumbel-Softmax 把 generator 的词表分布、Gumbel noise 和 temperature softmax 组合成可训练路径。}
{00:24:30--00:24:45}

训练时可以把 $y$ 当作 soft one-hot，再用 embedding matrix $E\in\mathbb{R}^{V\times d}$ 得到连续 embedding：
\[
e = E^\top y.
\]

这样 classifier 接收到的是所有词 embedding 的加权和，而不是单个离散词。虽然这不完全等同真实文本，但它允许梯度通过 $e$、$y$ 回到生成器参数 $\theta$。

\subsection{近似 one-hot 与梯度回传}

课程强调，Gumbel-Softmax 最终不是直接生成一个真实 token，而是用一个接近 one-hot 的分布近似 token。这个近似分布可以送入 embedding table，形成可微的“软 token”。

\framefig{0.90}{figures/fig08_gumbel_one_hot_approximation.jpg}
{使用 Gumbel-Softmax 分布近似 one-hot vector，再与 embedding table 相乘得到连续表示。}
{00:25:45--00:26:00}

因此，攻击训练可以形成完整反向路径：
\[
G_\theta(x)\rightarrow \pi
\rightarrow y_{\mathrm{gumbel}}
\rightarrow E^\top y_{\mathrm{gumbel}}
\rightarrow C(\cdot)
\rightarrow L_{\mathrm{adv}}.
\]

\framefig{0.90}{figures/fig09_gradient_through_autoencoder.jpg}
{由于 Gumbel-Softmax 提供连续近似，text classifier 的梯度可以回传穿过 auto-encoder。}
{00:27:15--00:27:30}

实际使用时常见两种变体。第一种是纯 soft relaxation：前向和反向都用连续 $y$。第二种是 straight-through estimator：前向用 hard one-hot token，反向时假装它是 soft $y$，让梯度穿过。前者训练更平滑，后者与真实离散输入更接近，但梯度估计偏差更明显。

\begin{importantbox}{Gumbel-Softmax 解决的不是“语义”问题}
    Gumbel-Softmax 只解决离散采样不可微的问题。它让生成器能从分类器 loss 得到梯度，但不自动保证生成句语义正确。语义保持仍要靠 reconstruction、similarity、POS、perplexity 或人工评估等约束。
\end{importantbox}

\subsection{本章小结}

文本生成式攻击的关键障碍是 token 采样不可微。Gumbel-Softmax 用 Gumbel noise 加 temperature softmax 近似 categorical sampling，使分类器梯度可以穿过 soft token 回到 generator。它把离散文本攻击暂时搬到连续空间中训练，再用近似 one-hot 的方式靠近真实 token。

% =====================================================================
\section{训练一个执行扰动的 Agent}
% =====================================================================

除了用 auto-encoder 直接生成整句，还可以把攻击看成一个 agent 的决策过程。Agent 读入句子，决定在哪些位置做什么扰动；分类器反馈当前扰动是否有效；训练目标是让 agent 学会用少量、合理的动作降低正确类别概率。

\framefig{0.90}{figures/fig10_perturbation_agent_goal.jpg}
{Perturbation agent 的目标：生成扰动样本，使 text classifier 对原类别的判断下降甚至翻转。}
{00:30:00--00:30:15}

\subsection{动作空间}

Agent 的动作可以设计成不同粒度。简单版本只在每个词上决定“保持不变”或“替换成某个候选词”；复杂版本可以允许替换、插入、删除、同义改写、形态变化等动作。课程中的表格展示的是对词进行替换类操作的动作集合。

\framefig{0.90}{figures/fig11_agent_action_space.jpg}
{Agent 的 action space 描述它能对文本做哪些扰动，例如保留当前词或替换成候选词。}
{00:31:45--00:32:00}

可以把决策写成策略：
\[
a_t\sim \pi_\theta(a_t\mid s_t),
\]
其中 $s_t$ 是当前句子状态，$a_t$ 是扰动动作。执行动作后得到新状态 $s_{t+1}$，也就是新的改写句子。

动作空间设计会影响攻击质量。动作太自由，容易产生不自然句子；动作太保守，攻击成功率可能很低。因此常见做法是先用同义词、词向量邻居或语言模型候选限制动作集合，再让 agent 在这些候选里选择。

\subsection{强化学习奖励}

因为动作是离散选择，强化学习是自然选择。课程用分类器对 ground truth class 的概率下降作为 reward：如果扰动后正确类别概率下降，agent 得到正奖励；如果概率没有下降或样本质量变差，则奖励较小或为负。

\framefig{0.90}{figures/fig12_reinforcement_learning_reward.jpg}
{强化学习奖励可以定义为真实类别概率的下降量，例如 $r=-\Delta p_y$。}
{00:33:45--00:34:00}

若当前句子为 $x_t$，执行动作后为 $x_{t+1}$，可以定义：
\[
r_t=p_C(y\mid x_t)-p_C(y\mid x_{t+1}).
\]

若 $r_t>0$，表示动作降低了模型对正确标签的信心；若 $r_t<0$，表示动作反而让模型更确定。对 targeted attack，也可以把 reward 改成目标类别概率的上升量：
\[
r_t=p_C(y_{\mathrm{tar}}\mid x_{t+1})-p_C(y_{\mathrm{tar}}\mid x_t).
\]

训练时可用 policy gradient：
\[
\nabla_\theta J(\theta)
=
\E_{\pi_\theta}
\left[
\sum_t R_t\nabla_\theta \log \pi_\theta(a_t\mid s_t)
\right],
\]
其中 $R_t$ 是从当前步骤开始的累积奖励。

\subsection{与 greedy 搜索的关系}

Perturbation agent 和 greedy/WIR 的目标相似：都是找一串离散替换动作。不同之处在于 greedy 是手写规则，每一步选当前最优；agent 是从经验中学策略，可以把多步效果纳入考虑。理论上，agent 可以学到“某个动作单看不强，但为后续动作铺路”的策略。

不过，强化学习训练也更不稳定，且需要大量 interaction。若分类器查询昂贵，或奖励设计不够精细，agent 可能学到低质量扰动。因此实际评估仍要看修改比例、语义相似度、语法自然度和查询成本。

\subsection{本章小结}

Perturbation agent 把文本攻击写成序列决策问题。动作空间决定 agent 能改什么，reward 决定 agent 追求什么。用真实类别概率下降作为奖励，可以训练 agent 生成更容易欺骗分类器的文本；但离散动作和样本质量约束仍是核心难题。

% =====================================================================
\section{防御一：Adversarial Training}
% =====================================================================

课程转入 defense 后，首先讲的是最直接的策略：在训练过程中加入 adversarial examples，让模型见过这些扰动模式。图像领域已有类似思想；在文本领域，困难在于 adversarial samples 的生成更贵，而且文本扰动必须保持语义。

\framefig{0.90}{figures/fig13_adversarial_training_current_model.jpg}
{Adversarial training 可以定期用当前模型生成 adversarial samples，再把它们加入训练。}
{00:35:45--00:36:00}

\subsection{标准训练流程}

标准 adversarial training 可写成：
\[
\min_\theta
\E_{(x,y)\sim\mathcal{D}}
\left[
\max_{x'\in \mathcal{A}(x)}
L(C_\theta(x'),y)
\right].
\]

其中 $\mathcal{A}(x)$ 是允许的扰动集合，包括修改比例、语义相似、POS consistency 等约束。内层最大化找最能让模型出错的扰动，外层最小化让模型在这些扰动上仍然正确。

\framefig{0.90}{figures/fig14_adversarial_training_loop.jpg}
{训练循环：攻击算法生成 adversarial training data，模型再用这些样本更新参数。}
{00:37:15--00:37:30}

实际训练中，不可能每一步都精确求解内层最大化，通常用已有攻击算法近似生成样本。例如每若干 epoch 用当前模型生成一批 adversarial examples，加入训练集，继续训练模型。这样模型会逐渐适应当前攻击器能找到的弱点。

\subsection{文本 adversarial training 的局限}

文本 adversarial training 的主要限制有三点：

\begin{enumerate}
    \item \textbf{成本高}：生成文本 adversarial examples 往往需要大量模型查询、候选过滤和语义检查。
    \item \textbf{攻击器依赖}：模型可能只对训练时使用的攻击方法鲁棒，对新攻击方法仍脆弱。
    \item \textbf{样本质量难控}：若加入低质量 adversarial samples，模型可能学到错误标签或偏离真实任务。
\end{enumerate}

因此，adversarial training 不是万能防御。它更像是一种压力训练：把已知攻击样本加入训练，让模型在这些局部区域变得更平滑、更稳定。

\subsection{本章小结}

Adversarial training 的思想简单但成本不低：先用攻击器生成难样本，再用这些样本训练鲁棒模型。对 NLP 来说，它必须建立在高质量攻击样本之上，否则防御可能只是记住低质量改写，而不是真正提升语义鲁棒性。

% =====================================================================
\section{防御二：Embedding 空间与 ASCC}
% =====================================================================

直接在离散文本上找 adversarial examples 很难，防御者也可以转向 embedding space。一个词的同义词往往在 embedding 空间中相近；若模型在这个邻域内都能保持预测稳定，就可能对同义替换攻击更鲁棒。

\subsection{Embedding ball adversarial training}

Embedding ball 方法把每个词的 embedding $e_i$ 看成连续向量，在其附近找扰动：
\[
e_i' = e_i + \delta_i,\qquad \|\delta_i\|\leq \epsilon.
\]

\framefig{0.90}{figures/fig15_embedding_ball_training.jpg}
{Embedding ball adversarial training 在词向量邻域中寻找扰动，并让模型对邻域内输入保持正确。}
{00:39:45--00:40:00}

训练目标类似：
\[
\min_\theta
\E_{(x,y)}
\left[
\max_{\|\delta_i\|\leq \epsilon}
L(C_\theta(e_1+\delta_1,\dots,e_n+\delta_n),y)
\right].
\]

优点是连续优化方便，可以用梯度方法找到 worst-case perturbation；缺点是 embedding ball 中的点未必对应真实词。模型在连续空间里变鲁棒，不一定等价于对真实同义词替换鲁棒。

\subsection{ASCC 的动机}

ASCC（Adversarial Sparse Convex Combination）试图解决这个问题。与其允许任意 $\epsilon$-ball 里的连续点，不如只考虑某个词的同义词集合构成的 convex hull。这样 adversarial embedding 仍是连续组合，但组合基底来自真实候选词。

\framefig{0.90}{figures/fig16_ascc_convex_hull_idea.jpg}
{ASCC 的基本想法：用同义词集合的 convex hull 作为扰动空间，而不是任意 embedding ball。}
{00:42:00--00:42:15}

设词 $w_i$ 的候选集合为：
\[
S_i=\{w_{i1},w_{i2},\dots,w_{ik}\},
\]
其 embedding 为 $e(w_{ij})$。ASCC 构造：
\[
\tilde{e}_i = \sum_{j=1}^{k}\alpha_{ij}e(w_{ij}),
\qquad
\alpha_{ij}\geq 0,\quad
\sum_{j=1}^{k}\alpha_{ij}=1.
\]

\framefig{0.90}{figures/fig17_ascc_linear_combination.jpg}
{ASCC 用候选同义词 embedding 的线性凸组合表示 adversarial embedding。}
{00:45:45--00:46:00}

这让扰动空间从“所有邻域点”变成“同义词候选张成的凸包”。如果候选集合质量较高，扰动就更接近可解释的词替换。

\subsection{稀疏性与最坏组合}

ASCC 还希望线性组合不要平均混合太多词，而是更接近少数候选词。因此目标中会加入鼓励 sparse coefficient 的项。一个概念化目标为：
\[
\max_{\alpha}
L(C_\theta(\tilde{e}_1,\dots,\tilde{e}_n),y)
-\lambda \sum_i H(\alpha_i),
\]
其中 $H(\alpha_i)$ 是位置 $i$ 上组合系数的熵。最大化时惩罚高熵分布，会使 $\alpha_i$ 更尖锐，接近选择少数候选词。

\framefig{0.90}{figures/fig18_ascc_sparse_coefficients.jpg}
{ASCC 在最大化分类损失的同时约束线性组合系数，使 adversarial embedding 更接近稀疏候选组合。}
{00:49:30--00:49:45}

随后模型在这些 worst-case convex combinations 上训练：
\[
\min_\theta
\E_{(x,y)}
\left[
L(C_\theta(\tilde{e}^{\star}),y)
\right],
\qquad
\tilde{e}^{\star}=\argmax_{\tilde{e}\in \mathrm{Conv}(S)}
L(C_\theta(\tilde{e}),y).
\]

ASCC 的优势是比任意 embedding ball 更贴近语义候选；局限是它高度依赖候选同义词集合。如果候选集合本身包含语义不等价的词，防御训练也会受到污染。

\subsection{用不鲁棒分类器做数据增强}

课程还提到一种数据增强思路：先训练一个不够鲁棒的分类器，用攻击算法从它身上生成 adversarial samples，再把这些样本加入训练数据，训练一个新的分类器。

\framefig{0.90}{figures/fig19_data_augmentation_unrobust_classifier.jpg}
{也可以先用不鲁棒分类器预生成 adversarial samples，再把它们加入训练集训练新模型。}
{00:52:15--00:52:30}

这种方法的优点是流程简单，且不需要在每次训练迭代中都运行昂贵攻击；缺点是生成样本只覆盖旧模型的脆弱点，新模型仍可能被其他攻击方式击中。

\begin{importantbox}{连续鲁棒性与离散鲁棒性}
    Embedding-space defense 让训练更容易优化，但文本最终仍是离散 token。若连续扰动空间和真实可替换词集合不匹配，模型可能在数学邻域里鲁棒，却在真实同义词攻击下仍不稳定。ASCC 的价值就在于把连续优化限制到同义词凸包附近。
\end{importantbox}

\subsection{本章小结}

Embedding ball adversarial training 把文本防御搬到连续向量空间中，方便使用梯度寻找最坏扰动；ASCC 进一步把扰动限制在同义词候选的 convex hull 内，让连续扰动更接近真实词替换。二者都试图提升模型对语义相近改写的稳定性，但都依赖候选集合和语义约束的质量。

% =====================================================================
\section{防御三：推理期检测与修复}
% =====================================================================

除了训练更鲁棒的模型，也可以在 inference time 检测输入是否被攻击，并尝试把它修复回原句。课程介绍的代表方法是 DISP，以及一个基于词频的简单检测/替换思路 FGWS。

\framefig{0.90}{figures/fig20_detecting_adversaries_outline.jpg}
{推理期防御关注 detecting adversaries during inference：输入进来后先判断是否含扰动，再尝试修复。}
{00:54:45--00:55:00}

\subsection{DISP：三段式检测与恢复}

DISP（Discriminate Perturbations）把问题拆成三个子模块：

\begin{enumerate}
    \item \textbf{Perturbation discriminator}：判断每个 token 是否被扰动。
    \item \textbf{Embedding estimator}：估计被扰动位置原本应该对应的 embedding。
    \item \textbf{Token recovery}：在 embedding corpus 中找最近词，把可疑 token 替换回去。
\end{enumerate}

\framefig{0.90}{figures/fig21_disp_three_modules.jpg}
{DISP 包含 perturbation discriminator、embedding estimator 和 token recovery 三个子模块。}
{00:56:00--00:56:15}

这种设计的直觉是：对抗样本中的替换词虽然在表面上像同义词，但在上下文中可能有异常痕迹。若能定位异常 token，再估计原始 token 的 embedding，就可以把句子修复到更接近 clean input 的状态。

\subsection{Embedding estimator}

Embedding estimator 的任务不是直接猜原词，而是先预测原词 embedding。它可以利用上下文表示、双向模型或邻近 token 信息，输出一个连续向量 $\hat{e}_i$：
\[
\hat{e}_i = R_\psi(h_i),
\]
其中 $h_i$ 是位置 $i$ 的上下文表示，$R_\psi$ 是 regression head。

\framefig{0.90}{figures/fig22_disp_embedding_estimator.jpg}
{Embedding estimator 利用上下文估计被扰动 token 原本的 embedding，再与词表 embedding 比较。}
{00:58:30--00:58:45}

随后 token recovery 可以执行最近邻检索：
\[
\hat{w}_i=\argmin_{w\in\mathcal{V}}
\|\hat{e}_i-e(w)\|_2.
\]

若 perturbation discriminator 认为某个位置没有被扰动，就不需要替换；若认为它被扰动，则用 $\hat{w}_i$ 替换当前 token。

\framefig{0.90}{figures/fig23_disp_regression_head.jpg}
{DISP 中的 regression head 负责把上下文表示映射回估计的原始 embedding。}
{01:01:00--01:01:15}

\subsection{训练与推理流程}

DISP 需要构造 clean sentence 与 perturbed sentence 的配对数据。训练时，系统知道哪些 token 被替换，因此可以监督 perturbation discriminator；同时也知道原词 embedding，因此可以训练 embedding estimator。

\framefig{0.90}{figures/fig24_disp_training_inference.jpg}
{DISP 训练时使用 clean/perturbed 配对监督各模块，推理时检测并恢复可疑扰动。}
{01:03:30--01:03:45}

推理时流程是：

\begin{enumerate}
    \item 输入一个可能被攻击的句子。
    \item discriminator 标出可疑 token。
    \item estimator 对可疑位置估计原始 embedding。
    \item recovery 从词表中找最近 token 进行替换。
    \item 将修复后的句子送入原分类器。
\end{enumerate}

DISP 的优点是能在模型预测前做预处理；局限是它假设扰动具有可检测模式。如果攻击者知道防御器存在，可能生成更像自然输入的替换，绕过 detector。

\subsection{FGWS：用词频变化发现可疑替换}

FGWS（Frequency-Guided Word Substitution）利用一个经验观察：许多文本攻击会把常见词替换成语义相近但频率较低的词，从而骗过模型。于是可以尝试把低频可疑词替换成更高频的对应词，再比较模型输出变化。

\framefig{0.90}{figures/fig25_fgws_probability_threshold.jpg}
{FGWS 比较原输入与替换后输入在模型预测概率上的差异；若差异超过阈值，就标记为 adversarial。}
{01:07:30--01:07:45}

一个简化检测规则为：
\[
\Delta p =
p_C(\hat{y}\mid \mathrm{swap}(x))
-p_C(\hat{y}\mid x),
\]
其中 $\hat{y}=C(x)$ 是原输入的预测类别，$\mathrm{swap}(x)$ 表示把低频词替换成高频候选后的句子。若：
\[
\Delta p>\gamma,
\]
则认为原输入可能是 adversarial input。

FGWS 的优点是简单，不需要复杂训练；缺点是启发式很强。它依赖“攻击倾向使用低频词”这个假设，对使用高频自然替换、上下文语言模型替换或语义更隐蔽的攻击可能不够有效。

\begin{warningbox}{检测式防御的共同风险}
    检测器本身也会成为攻击目标。若攻击者知道检测规则，就可以优化样本同时欺骗分类器和检测器。因此推理期检测适合作为多层防御的一部分，不应被视为单独充分的安全保证。
\end{warningbox}

\subsection{本章小结}

推理期防御不直接改变分类器，而是在输入进入分类器前检测和修复。DISP 通过 token-level detector、embedding estimator 和 nearest-neighbor recovery 还原可疑扰动；FGWS 通过低频词替换与概率差阈值做启发式检测。它们能补充 adversarial training，但都面临适应性攻击的挑战。

% =====================================================================
\section{总结与延伸}
% =====================================================================

P42 把 NLP evasion attack 的后半部分补完整：前半讲关注如何生成攻击样本，后半讲关注如何防御或检测攻击样本。整讲的主线是“离散文本如何被连续优化处理”。

在攻击侧，auto-encoder attack 把 adversarial example generation 写成训练目标：重建 loss 保持句子接近，similarity loss 保持语义接近，adversarial loss 推动分类器犯错。由于 token sampling 不可微，Gumbel-Softmax 用 soft one-hot 近似离散采样，让分类器梯度可以穿过 generator。另一条路线是训练 perturbation agent，把文本改写看成离散动作序列，用真实类别概率下降作为 reward。

在防御侧，adversarial training 让模型见过攻击样本；embedding-space training 把离散替换放到连续空间中寻找最坏扰动；ASCC 则进一步把扰动限制在同义词 convex hull 中，试图兼顾可优化性和语义合理性。推理期方法如 DISP 和 FGWS 不重新训练主模型，而是在输入侧检测、修复或替换可疑 token。

这几类方法共同说明，NLP 对抗攻防的本质不是单个算法名，而是三个问题的组合：

\begin{itemize}
    \item \textbf{优化问题}：离散 token 如何接受梯度或搜索？
    \item \textbf{语义问题}：改写后是否仍是同一个输入？
    \item \textbf{安全问题}：防御是否只挡住已知攻击，还是能应对自适应攻击者？
\end{itemize}

从实践角度看，单一防御很难覆盖所有攻击。较可靠的系统通常需要组合：训练时使用高质量 adversarial examples 增强鲁棒性，模型设计时检查 embedding 与同义词邻域的稳定性，推理时加入检测和异常监控，评估时用多种攻击方法与人工语义检查验证。NLP 模型越接近真实应用，越不能只看 clean accuracy；它还必须经受“语义不变但表面改变”的压力测试。

\end{document}
